% GENERATED FILE. Edit canonical lessons, then run npm run generate:books.
% canonical-source-hash: fnv1a64:2278d7802afaa13c
% canonical-lessons: FR-C158-donc, FR-C158-pourtant, FR-C158-de-toute-facon, FR-C158-au-contraire, FR-C158-bien-que

\chapter{So, However, Anyway, On the contrary, Although}
\label{ch:fr-so-however-anyway-on-the-contrary-although}
\hlchaptermodality{158}

\begin{chapteropening}
\textbf{By the end of this chapter:} \emph{I can say so, however, anyway, on the contrary and although.}

Complete the last lesson of chapter 158: I can say so, however, anyway, on the contrary and although.
\end{chapteropening}

\section[donc]{donc — so, therefore}

\label{lesson:FR-C158-donc}



\begin{quote}
Before the new one: say the French for behind, then the French for beside.
\end{quote}

\subsection*{You'll want to know: donc}
\textbf{donc} — \textquotedblleft{}so, therefore\textquotedblright{}.

\begin{grammarlens}[title={What you've built}]
One more word for this chapter.
\end{grammarlens}

\subsection*{Your turn}
\begin{itemize}
\raggedright
  \item \emph{Say it:} \emph{donc}
  \item \emph{Say it:} \emph{donc}, once more
  \item \emph{Say it:} say \emph{donc}
  \item \emph{Recall:} say the French for behind, then the French for beside, then say \emph{donc} again
\end{itemize}

\begin{tcolorbox}[breakable,colback=teal!4,colframe=teal!35!black,title={Before you move on}]
What does donc mean? (\textquotedblleft{}So, therefore\textquotedblright{}.) Say it once more.
\end{tcolorbox}
\section[pourtant]{pourtant — however, yet}

\label{lesson:FR-C158-pourtant}



\begin{quote}
Before the new one: say the French for beside, then the French for so.
\end{quote}

\subsection*{You'll want to know: pourtant}
\textbf{pourtant} — \textquotedblleft{}however, yet\textquotedblright{}.

\begin{grammarlens}[title={What you've built}]
One more word for this chapter.
\end{grammarlens}

\subsection*{Your turn}
\begin{itemize}
\raggedright
  \item \emph{Say it:} \emph{pourtant}
  \item \emph{Say it:} \emph{pourtant}, once more
  \item \emph{Say it:} say \emph{pourtant}
  \item \emph{Recall:} say the French for beside, then the French for so, then say \emph{pourtant} again
\end{itemize}

\begin{tcolorbox}[breakable,colback=teal!4,colframe=teal!35!black,title={Before you move on}]
What does pourtant mean? (\textquotedblleft{}However, yet\textquotedblright{}.) Say it once more.
\end{tcolorbox}
\section[de toute façon]{de toute façon — anyway}

\label{lesson:FR-C158-de-toute-facon}



\begin{quote}
Before the new one: say the French for so, then the French for however.
\end{quote}

\subsection*{You'll want to know: de toute façon}
\textbf{de toute façon} — \textquotedblleft{}anyway\textquotedblright{}.

\begin{grammarlens}[title={What you've built}]
One more word for this chapter.
\end{grammarlens}

\subsection*{Your turn}
\begin{itemize}
\raggedright
  \item \emph{Say it:} \emph{de toute façon}
  \item \emph{Say it:} \emph{de toute façon}, once more
  \item \emph{Say it:} say \emph{de toute façon}
  \item \emph{Recall:} say the French for so, then the French for however, then say \emph{de toute façon} again
\end{itemize}

\begin{tcolorbox}[breakable,colback=teal!4,colframe=teal!35!black,title={Before you move on}]
What does de toute façon mean? (\textquotedblleft{}Anyway\textquotedblright{}.) Say it once more.
\end{tcolorbox}
\section[au contraire]{au contraire — on the contrary}

\label{lesson:FR-C158-au-contraire}



\begin{quote}
Before the new one: say the French for however, then the French for anyway.
\end{quote}

\subsection*{You'll want to know: au contraire}
\textbf{au contraire} — \textquotedblleft{}on the contrary\textquotedblright{}.

\begin{grammarlens}[title={What you've built}]
One more word for this chapter.
\end{grammarlens}

\subsection*{Your turn}
\begin{itemize}
\raggedright
  \item \emph{Say it:} \emph{au contraire}
  \item \emph{Say it:} \emph{au contraire}, once more
  \item \emph{Say it:} say \emph{au contraire}
  \item \emph{Recall:} say the French for however, then the French for anyway, then say \emph{au contraire} again
\end{itemize}

\begin{tcolorbox}[breakable,colback=teal!4,colframe=teal!35!black,title={Before you move on}]
What does au contraire mean? (\textquotedblleft{}On the contrary\textquotedblright{}.) Say it once more.
\end{tcolorbox}
\section[bien que]{bien que — although}

\label{lesson:FR-C158-bien-que}



\begin{quote}
Before the new one: say the French for anyway, then the French for on the contrary.
\end{quote}

\subsection*{You'll want to know: bien que}
\textbf{bien que} — \textquotedblleft{}although\textquotedblright{}.

\begin{grammarlens}[title={What you've built}]
One more word for this chapter.
\end{grammarlens}

\subsection*{Your turn}
\begin{itemize}
\raggedright
  \item \emph{Say it:} \emph{bien que}
  \item \emph{Say it:} \emph{bien que}, once more
  \item \emph{Say it:} say \emph{bien que}
  \item \emph{Recall:} say the French for anyway, then the French for on the contrary, then say \emph{bien que} again
\end{itemize}

\begin{tcolorbox}[breakable,colback=teal!4,colframe=teal!35!black,title={Before you move on}]
What does bien que mean? (\textquotedblleft{}Although\textquotedblright{}.) Say it once more.
\end{tcolorbox}
